\section*{Sets and Indices}

\begin{itemize}
    \item Products \(P = \{\text{Meaties}, \text{Yummies}\}\), indexed by \(p\).
    \item Let \(m\) denote product Meaties and \(y\) denote product Yummies, matching the code aliases \texttt{m} and \texttt{y}.
\end{itemize}

\section*{Parameters}

\begin{itemize}
    \item Selling price per pack of product \(p\): \(s_p\).
    \item Grains required per pack of product \(p\): \(g_p\).
    \item Meat required per pack of product \(p\): \(r_p\).
    \item Variable production cost per pack of product \(p\): \(v_p\).
    \item Production capacity for product \(p\): \(u_p\), when provided in the data; if blank, no capacity constraint is added in the code.
    \item Maximum available grains: \(G\) (code source: \texttt{max\_grains}).
    \item Cost per pound of grains: \(c^g\) (code source: \texttt{price\_grains}).
    \item Maximum available meat: \(M\) (code source: \texttt{max\_meat}).
    \item Cost per pound of meat: \(c^r\) (code source: \texttt{price\_meat}).
\end{itemize}

For this instance, the data are:
\[
G = 400000,\qquad c^g = 0.20,\qquad M = 300000,\qquad c^r = 0.50.
\]

Product-specific values:
\[
\begin{array}{c|ccccc}
p & s_p & g_p & r_p & v_p & u_p \\
\hline
m\;(\text{Meaties}) & 2.80 & 2.0 & 3.0 & 0.25 & 90000 \\
y\;(\text{Yummies}) & 2.00 & 3.0 & 1.5 & 0.20 & \text{none}
\end{array}
\]

The code computes per-pack profit coefficients as
\[
\pi_p = s_p - v_p - (g_p c^g + r_p c^r).
\]
Hence
\[
\pi_m = 2.80 - 0.25 - (2.0\cdot 0.20 + 3.0\cdot 0.50) = 0.65,
\]
\[
\pi_y = 2.00 - 0.20 - (3.0\cdot 0.20 + 1.5\cdot 0.50) = 0.45.
\]

\section*{Decision Variables}

\begin{itemize}
    \item Number of Meaties packs produced: \(x_m\) (code: \texttt{x\_meaties}).
    \item Number of Yummies packs produced: \(x_y\) (code: \texttt{x\_yummies}).
\end{itemize}

Both variables are continuous and nonnegative:
\[
x_m \ge 0,\qquad x_y \ge 0.
\]

\section*{Objective}

Maximize total profit:
\[
\max \; \pi_m x_m + \pi_y x_y.
\]

Substituting the coefficients used by the code:
\[
\max \; 0.65\,x_m + 0.45\,x_y.
\]

\section*{Constraints}

Raw material availability:
\[
g_m x_m + g_y x_y \le G,
\]
\[
r_m x_m + r_y x_y \le M.
\]

For this instance:
\[
2.0\,x_m + 3.0\,x_y \le 400000,
\]
\[
3.0\,x_m + 1.5\,x_y \le 300000.
\]

Production capacity constraints are added only when the corresponding CSV field is nonempty. Therefore:
\[
x_m \le u_m = 90000.
\]

No production-capacity constraint is added for Yummies because its \texttt{Production\_capacity} field is blank in the data.

\section*{Notes on Implementation}

\begin{itemize}
    \item The implementation builds a linear program in PuLP-compatible syntax and solves it with \texttt{GUROBI\_CMD(msg=0)}.
    \item The objective coefficients are not read directly from a profit column; instead, the code explicitly computes profit per pack as selling price minus variable cost minus raw material costs:
    \[
    \pi_p = s_p - v_p - (g_p c^g + r_p c^r).
    \]
    \item Decision variables are continuous nonnegative variables; no integrality restrictions are imposed on pack counts.
    \item Capacity constraints are conditional in the code. A product receives an upper-bound constraint only if its \texttt{Production\_capacity} entry is present and nonblank in the CSV input.
\end{itemize}
